\par\SourceLeaf{30}
{\rightskip=0pt plus 1em What is distinctive about the principle of virtual velocities is that it provides a general formula for solving all problems of statics and thus stands in for all the other principles. Yet it does not show its credentials for doing so directly enough to make itself seem plausible as soon as it is stated.\par}

In \textit{this} respect, the principle I am about to put forward seems to have the advantage. It also has a second advantage, in that it encompasses the law of motion and of rest in exactly the same way and with the utmost generality. In the gradual development of science and in the instruction of the individual, it is entirely proper for what is easier to precede what is harder, for what is simpler to precede what is more involved, and for the particular to precede the general. Once the mind has reached the higher standpoint, however, it demands the reverse course, in which all of statics appears as only a very special case of mechanics. Even the geometer mentioned above seems to value this. He regards it as an advantage of the principle of least action that it encompasses equilibrium and motion together, provided equilibrium is expressed so that the quantities of vis viva are minima in both cases. This remark seems more witty than true, since the minimum occurs in an entirely different respect in each case.

The new principle is as follows.

\textit{A system of material points connected to one another in any manner whatsoever and whose motions are also subject to any external restrictions moves at every instant in the closest possible agreement with free motion, or under the least possible constraint. Here the constraint on the entire system in each interval of time is measured by the sum of the products of the square of each point’s deviation from its free motion and the mass of that point.}

Let $m,m',m''$ and so on be the masses of the points, $a,a',a''$ and so on their positions at time $t$, and $b,b',b''$ and so on the positions they would occupy after the infinitesimal
